%-----------------------------------------------------------------------------------------------------------------------------------------------%
%	The MIT License (MIT)
%
%	Copyright (c) 2015 Jan Küster
%
%	Permission is hereby granted, free of charge, to any person obtaining a copy
%	of this software and associated documentation files (the "Software"), to deal
%	in the Software without restriction, including without limitation the rights
%	to use, copy, modify, merge, publish, distribute, sublicense, and/or sell
%	copies of the Software, and to permit persons to whom the Software is
%	furnished to do so, subject to the following conditions:
%	
%	THE SOFTWARE IS PROVIDED "AS IS", WITHOUT WARRANTY OF ANY KIND, EXPRESS OR
%	IMPLIED, INCLUDING BUT NOT LIMITED TO THE WARRANTIES OF MERCHANTABILITY,
%	FITNESS FOR A PARTICULAR PURPOSE AND NONINFRINGEMENT. IN NO EVENT SHALL THE
%	AUTHORS OR COPYRIGHT HOLDERS BE LIABLE FOR ANY CLAIM, DAMAGES OR OTHER
%	LIABILITY, WHETHER IN AN ACTION OF CONTRACT, TORT OR OTHERWISE, ARISING FROM,
%	OUT OF OR IN CONNECTION WITH THE SOFTWARE OR THE USE OR OTHER DEALINGS IN
%	THE SOFTWARE.
%	
%
%-----------------------------------------------------------------------------------------------------------------------------------------------%


%============================================================================%
%
%	DOCUMENT DEFINITION
%
%============================================================================%

%we use article class because we want to fully customize the page and dont use a cv template
\documentclass[10pt,A4]{article}	


%----------------------------------------------------------------------------------------
%	ENCODING
%----------------------------------------------------------------------------------------

%we use utf8 since we want to build from any machine
\usepackage[utf8]{inputenc}		

%----------------------------------------------------------------------------------------
%	LOGIC
%----------------------------------------------------------------------------------------

% provides \isempty test
\usepackage{xifthen}

%----------------------------------------------------------------------------------------
%	FONT
%----------------------------------------------------------------------------------------

% some tex-live fonts - choose your own

%\usepackage[defaultsans]{droidsans}
%\usepackage[default]{comfortaa}
%\usepackage{cmbright}
\usepackage[default]{raleway}
%\usepackage{fetamont}
%\usepackage[default]{gillius}
%\usepackage[light,math]{iwona}
%\usepackage[thin]{roboto} 

% set font default
\renewcommand*\familydefault{\sfdefault} 	
\usepackage[T1]{fontenc}

% more font size definitions
\usepackage{moresize}		

% hyperlink definition
\usepackage[hidelinks]{hyperref}


%----------------------------------------------------------------------------------------
%	PAGE LAYOUT  DEFINITIONS
%----------------------------------------------------------------------------------------

%debug page outer frames
%\usepackage{showframe}			


%define page styles using geometry
\usepackage[a4paper]{geometry}		

% for example, change the margins to 2 inches all round
\geometry{top=1.75cm, bottom=-.6cm, left=1.5cm, right=1.5cm} 	

%use customized header
\usepackage{fancyhdr}				
\pagestyle{fancy}

%less space between header and content
\setlength{\headheight}{-5pt}		

% package for multiple optional arguments
\usepackage{xparse}

%customize entries left, center and right
\lhead{}
\chead{\small{$\cdot$ ORCID: 0000-0001-9152-8094  $\cdot$  
	\textcolor{sectcol}{\textbf{\href{mailto:nikan.doosti@outlook.com}{nikan.doosti@outlook.com}}}  $\cdot$ 
	\textcolor{sectcol}{\textbf{\href{https://nikronic.github.io}{nikronic.github.io}}} $\cdot$
	}}
\rhead{}


%indentation is zero
\setlength{\parindent}{0mm}

%----------------------------------------------------------------------------------------
%	TABLE /ARRAY DEFINITIONS
%---------------------------------------------------------------------------------------- 

%for layouting tables
\usepackage{multicol}			
\usepackage{multirow}

% for aligning text (non-math)
\usepackage{amsmath}

%extended aligning of tabular cells
\usepackage{array}

\newcolumntype{x}[1]{%
>{\raggedleft\hspace{0pt}}p{#1}}%


%----------------------------------------------------------------------------------------
%	GRAPHICS DEFINITIONS
%---------------------------------------------------------------------------------------- 

%for header image
\usepackage{graphicx}

%for floating figures
\usepackage{wrapfig}
\usepackage{float}
%\floatstyle{boxed} 
%\restylefloat{figure}

%for drawing graphics		
\usepackage{tikz}				
\usetikzlibrary{shapes, backgrounds,mindmap, trees}


%----------------------------------------------------------------------------------------
%	Color DEFINITIONS
%---------------------------------------------------------------------------------------- 

\usepackage{color}

%accent color
\definecolor{sectcol}{RGB}{51,204,204}


%dark background color
\definecolor{bgcol}{RGB}{110,110,110}

%light background / accent color
\definecolor{softcol}{RGB}{225,225,225}


%============================================================================%
%
%
%	DEFINITIONS
%
%
%============================================================================%

%----------------------------------------------------------------------------------------
% 	HEADER
%----------------------------------------------------------------------------------------

% remove top header line
\renewcommand{\headrulewidth}{0pt} 
%
%remove botttom header line
\renewcommand{\footrulewidth}{0pt}	  	
%
%remove pagenum
\renewcommand{\thepage}{}	

%remove section num		
\renewcommand{\thesection}{}			

%----------------------------------------------------------------------------------------
% 	ARROW GRAPHICS in Tikz
%----------------------------------------------------------------------------------------

% a six pointed arrow poiting to the left
\newcommand{\tzlarrow}{(0,0) -- (0.2,0) -- (0.3,0.2) -- (0.2,0.4) -- (0,0.4) -- (0.1,0.2) -- cycle;}	

% include the left arrow into a tikz picture
% param1: fill color
%
\newcommand{\larrow}[1]
{\begin{tikzpicture}[scale=0.58]
	 \filldraw[fill=#1!100,draw=#1!100!black]  \tzlarrow
 \end{tikzpicture}
}

% a six pointed arrow poiting to the right
\newcommand{\tzrarrow}{ (0,0.2) -- (0.1,0) -- (0.3,0) -- (0.2,0.2) -- (0.3,0.4) -- (0.1,0.4) -- cycle;}

% include the right arrow into a tikz picture
% param1: fill color
%
\newcommand{\rarrow}
{\begin{tikzpicture}[scale=0.7]
	\filldraw[fill=sectcol!100,draw=sectcol!100!black] \tzrarrow
 \end{tikzpicture}
}



%----------------------------------------------------------------------------------------
%	custom sections
%----------------------------------------------------------------------------------------

% create a coloured box with arrow and title as cv section headline
% param 1: section title
%
\newcommand{\cvsection}[1]
{
\colorbox{sectcol}{\mystrut \makebox[1\linewidth][l]{
\larrow{bgcol} \hspace{-8pt} \larrow{bgcol} \hspace{-8pt} \larrow{bgcol} \textcolor{white}{\textbf{#1}}\hspace{4pt}
}}\\
}

%create a coloured arrow with title as cv meta section section
% param 1: meta section title
%
\newcommand{\metasection}[2]
{
\begin{tabular*}{1\textwidth}{p{2.0cm} p{10cm}}
\larrow{sectcol}	\normalsize{\textbf{#1}}&#2\\[12pt]
\end{tabular*}
}

%----------------------------------------------------------------------------------------
%	 CV EVENT
%----------------------------------------------------------------------------------------

% creates a stretched box as cv entry headline followed by two paragraphs about 
% the work you did
% param 1:	event time i.e. 2014 or 2011-2014 etc.
% param 2:	event name (what did you do?)
% param 3:	institution (where did you work / study)
% param 4:	description/summary
% param 5:	some words about your contributions in list form
%
\newcolumntype{P}[1]{>{\centering\arraybackslash}p{#1}}
\newcommand{\cvevent}[4]
{
	\vspace{8pt}

    \begin{tabular*}{1\textwidth}{P{4.6cm} P{5.7cm} P{6.6cm}}
        \textcolor{bgcol}{#1}& \textbf{#2} & 
        {\textbf{#3}}
    \end{tabular*}
    
	\vspace{-9pt}
	\textcolor{softcol}{\hrule}
	\vspace{6pt}
	{#4}
	\vspace{6pt}
}

% creates a list with left-arrows for show-casing the works done in a "cvevent" command
% 
\newcommand{\cveventitem}[1]
{
	\item[] \larrow{sectcol}  #1
}

% creates a stretched box as 
\newcommand{\cveventmeta}[2]
{
	\mbox{\mystrut \hspace{87pt}\textit{#1}}\\
	#2
}

%----------------------------------------------------------------------------------------
% CUSTOM STRUT FOR EMPTY BOXES
%----------------------------------------- -----------------------------------------------
\newcommand{\mystrut}{\rule[-.3\baselineskip]{0pt}{\baselineskip}}

%----------------------------------------------------------------------------------------
% CUSTOM LOREM IPSUM
%----------------------------------------------------------------------------------------
\newcommand{\lorem}
{Lorem ipsum dolor sit amet, consectetur adipiscing elit. Donec a diam lectus.}



%============================================================================%
%
%
%
%	DOCUMENT CONTENT
%
%
%
%============================================================================%
\begin{document}


%use our custom fancy header definitions
\pagestyle{fancy}	


%---------------------------------------------------------------------------------------
%	TITLE HEADLINE
%----------------------------------------------------------------------------------------
\vspace{-20.55pt}

% use this for multiple words like working titles etc.
%\hspace{-0.25\linewidth}\colorbox{bgcol}{\makebox[1.5\linewidth][c]{\hspace{46pt}\HUGE{\textcolor{white}{\textsc{Jan Küster}} } \textcolor{sectcol}{\rule[-1mm]{1mm}{0.9cm}} \parbox[b]{5cm}{   \large{ \textcolor{white}{{IT Consultant}}}\\
% \large{ \textcolor{white}{{Resume}}}}
%}}

% use this for single words, e.g. CV or RESUME etc.
\hspace{-0.25\linewidth}\colorbox{bgcol}{\makebox[1.5\linewidth][c]{\HUGE{\textcolor{white}{\textsc{Nikan Doosti}} }\textcolor{sectcol}{\rule[-1mm]{1mm}{0.9cm}} \HUGE{\textcolor{white}{\textsc{Resume}} } }}


%----------------------------------------------------------------------------------------
%	HEADER IMAGE
%----------------------------------------------------------------------------------------

\begin{figure}[H]
\begin{flushright}
	\includegraphics[width=0.2\linewidth]{qr-github.png}
\end{flushright}
\end{figure}

%---------------------------------------------------------------------------------------
%	QR CODE (optional)
%----------------------------------------------------------------------------------------
%\vspace{-136pt}
%\hspace{0.75\linewidth}
%\includegraphics[width=103pt]{qrcode}
%\normalsize
%\vspace{88pt}

%---------------------------------------------------------------------------------------
%	META SECTION
%----------------------------------------------------------------------------------------

\vspace{-114pt}

\metasection{Status:}{M.Sc. Computer Engineering, Artificial Intelligence}
\metasection{Tech:}{PyTorch, Python, Scipy/Numpy}
\metasection{Fields:}{Deep Learning, Computational Fabrication, Visual Computing, Computer Graphics}
\vspace{12pt}  % 6pt+12pt (from line space of metasection command)

%---------------------------------------------------------------------------------------
%	SUMMARAY (optional)
%----------------------------------------------------------------------------------------

\cvsection{Summary}\\
I am a final-year master student in computer engineering at IUST focusing on \emph{applications of artificial intelligence in science}. Before that, I was mainly focused on computer vision particularly \emph{image processing} and applications of machine learning in technology. Currently, I am looking for a PhD position in areas aforementioned above, and generally \emph{I love to do interdisiplineary work} between AI and science.\\[-2pt]

%============================================================================%
%
%	CV SECTIONS AND EVENTS (MAIN CONTENT)
%
%============================================================================%

%---------------------------------------------------------------------------------------
%	EXPERIENCE
%----------------------------------------------------------------------------------------
\cvsection{Experience}


%\textcolor{softcol}{\hrule}

%
\cvevent{Jul, 2020 - Mar, 2021}{Research Assistant (Remote)}{Max Planck Institute for Informatics}{This was my first mutlidisciplinary work which led to publication of my first-ever academic paper at ACM Symposium on Computational Fabrication (SCF21)!}

My responsibilities:
\begin{itemize}
	\itemsep -0.2em
	\cveventitem{Studying numerical simulation concepts such as FEM for solving Topology Optimization (TO) problem}
	\cveventitem{Studying and evaluating recent papers in implicit neural networks and solving PDEs using deep learning}
	\cveventitem{Implementing and reproducing results of the main references}
	\cveventitem{Developing a new system to solve TO in fully self-supervised manner via neural fields}
\end{itemize}
Supervisor: \href{https://cam.mpi-inf.mpg.de/?view=people_vahid}{\underline{Dr. Vahid Babaei}} (Computational manufacturing researcher at Max Planck Institute for Informatics and group leader at \href{https://cam.mpi-inf.mpg.de/?view=home}{\underline{CAM} group})
\vspace{12pt}

%
\cvevent{Aug, 2016 - Jun, 2020}{Freelancer}{Independent}{To help with the expenses and also learn more, during my study in BSc and earlier stages of MSc, I was working as a freelancer partially.}

My responsibilities:
\begin{itemize}
	\itemsep -0.2em
	\cveventitem{Writing small pieces of software for desktop computers mostly in C\#}
	\cveventitem{Designing, consulting and implementing standard algorithms and data structures for designing efficient systems}
	\cveventitem{Designing, consulting and implementing evolutionary algorithms such as genetics for developing efficient systems in different topics}
	\cveventitem{Designing, consulting and implementing classical computer vision and deep learning based models for different topics from reading car plates, tracking speed of cars, counting objects and measuring surfaces.}
\end{itemize}
\vspace{12pt}



%---------------------------------------------------------------------------------------
%	PUBLICATION SECTION
%--------------------------------------------------------------------------------------
\cvsection{Publications}

\begin{itemize}
	\cveventitem{Doosti, Nikan, Julian Panetta, and Vahid Babaei. "Topology Optimization via Frequency Tuning of Neural Design Representations." In \emph{Symposium on Computational Fabrication}, pp. 1-9. 2021 (\href{https://dl.acm.org/doi/abs/10.1145/3485114.3485124}{\underline{Abstract/PDF})}}
\end{itemize}
\vspace{12pt}



%---------------------------------------------------------------------------------------
%	EDUCATION SECTION
%--------------------------------------------------------------------------------------
\cvsection{Education}

\cvevent{Aug, 2015 - Aug, 2019}{BSc Computer Engineering}{University of Guilan}{}

Awards:
\begin{itemize}
	\itemsep -0.2em
	\cveventitem{Full tuition fee waiver}
	\cveventitem{Overall grade: 18.63/20}
% end of page, use "artificial footer"

%-------------------------------------------------------------------------------------------------
%	ARTIFICIAL FOOTER (fancy footer cannot exceed linewidth) 
%--------------------------------------------------------------------------------------------------
\null
\vspace*{\fill}
\hspace{-0.25\linewidth}\colorbox{bgcol}{\makebox[1.5\linewidth][c]{\mystrut \small \textcolor{white}{\href{https://nikronic.github.io}{nikronic.github.io}} $\cdot$ \textcolor{white}{\href{https://github.com/Nikronic}{github.com/Nikronic}} $\cdot$ \textcolor{white}{updated: 12.12.2021}}}


% start of new page, end of "artificial footer" - cont. section education
	\cveventitem{Ranked 3rd}
	\cveventitem{Exceptional talent at Department of Computer Engineering}
\end{itemize}

Related coursework:
\begin{itemize}
	\itemsep -0.2em
	\cveventitem{Algorithms and data structure}
	\cveventitem{Artificial Intelligence and Computational Intelligence}
	\cveventitem{Classical computer vision, statistical NLP}
	\cveventitem{Signal and systems}
	\cveventitem{Digital electronics and computer aided design (VHDL)}
\end{itemize}

Teacher assistant experience::
\begin{itemize}
	\itemsep -0.2em
	\cveventitem{Advanced programming (Head TA), instructor: \emph{Dr. Mirroshandel}, Fall 2018}
	\cveventitem{Algorithms, instructor: \emph{Dr. Shakeri}, Spring 2017}
	\cveventitem{Algorithms (Head TA), instructor: \emph{Dr. Shakeri}, Fall 2018}
	\cveventitem{Computational Intelligence (Head TA), instructor: \emph{Dr. Shakeri}, Spring 2018}
\end{itemize}
My primary responsibilities as the TA were assessing and designing assignments and teaching as the head TA.
\vspace{12pt}


%
\cvevent{Aug, 2019 - Present}{MSc Computer Engineering}{Iran University of Sc. and Tech. (IUST)}{}

Awards:
\begin{itemize}
	\itemsep -0.2em
	\cveventitem{Full tuition fee waiver}
	\cveventitem{Overall grades (so far): 17.17/20}
	\cveventitem{Accepted as an exceptional talent (no entry exam)}
\end{itemize}
Thesis: Toward super-resolution neural topology optimization (in progress)\\
Supervisor: \href{http://webpages.iust.ac.ir/mozayani/}{\underline{Dr. Nasser Mozayani}} (Associate professor of department of computer engineering)
\vspace{12pt}



%---------------------------------------------------------------------------------------
%	ACTIVITIES SECTION
%--------------------------------------------------------------------------------------
\cvsection{Voluntary Activities}

%
\cvevent{Oct, 2018 - Present}{Lecturer}{Rasht School of AI (\href{http://schoolofai.ir}{schoolofai.ir})}{School of AI is a community of passionate people about learning and teaching AI, mostly enabled by students from different fields of study. In this community, we would like to personalize learning paths and create a network to learn more about machine learning and deep learning and their applications. All sessions are free and recordings are available publicly.\\
You can find the material here: \href{https://github.com/rasht-school-of-ai}{\underline{github.com/rasht-school-of-ai}}}
\vspace{12pt}


%
\cvevent{Sep, 2020- Aug, 2021}{Maintainer and Developer}{IUST Projects (\href{http://iust-projects.ir}{iust-projects.ir})}{IUST Projects is an open GitHub organization with a focus on showcasing and maintaining projects created at Iran University of Science and Technology. This website and its \emph{organization is maintained by its true owners, the students} and the contributors outside of the university, so no, this is not an official university website and it is open to anybody who wants to join!\\
You can find all the materials here: \href{https://github.com/iust-projects}{\underline{github.com/iust-projects}}}

Main works:
\begin{itemize}
	\itemsep -0.2em
	\cveventitem{Assignments and solutions in form of blog posts for better readability}
	\cveventitem{Hosts course notes, curriculum, etc of AI major}
\end{itemize}
\vspace{12pt}


%
\cvevent{Nov, 2018 - Present}{Member}{PyTorch Discuss Forum}{Since the day I started deep learning with PyTorch, I started reading topics around my questions and tried to answer questions that I was not familiar with which helped me to dive deeper into the framework, learn many of hacks and return the favor! At the time of writing, \emph{I've visited the forum for 851 days, 563 posted replies, 10000 read posts and 184 solutions (All-time 15th)}\\
View the summary: \href{https://discuss.pytorch.org/u/nikronic/summary}{\underline{discuss.pytorch.org/u/nikronic/summary}}}
\vspace{12pt}
\vspace{18pt}

%-------------------------------------------------------------------------------------------------
%	ARTIFICIAL FOOTER (fancy footer cannot exceed linewidth) 
%--------------------------------------------------------------------------------------------------

\null
\vspace*{\fill}
\hspace{-0.25\linewidth}\colorbox{bgcol}{\makebox[1.5\linewidth][c]{\mystrut \small \textcolor{white}{\href{https://nikronic.github.io}{nikronic.github.io}} $\cdot$ \textcolor{white}{\href{https://github.com/Nikronic}{github.com/Nikronic}} $\cdot$ \textcolor{white}{updated: 12.12.2021}}}


%---------------------------------------------------------------------------------------
%	Selected Projects SECTION
%--------------------------------------------------------------------------------------
\cvsection{Selected Projects}

%
\cvevent{2016}{Paper Implementation}{Open Source}{Implementation of \emph{Using Genetic Algorithms for Multi-depot Vehicle Routing (Ombuki-Berman et al. 2009)} using Python as part of \emph{Computational Intelligence} course.\\
You can find the source code here: \href{https://github.com/Nikronic/Optimized-MDVRP}{\underline{github.com/Nikronic/Optimized-MDVRP}}}

Main contributions:
\begin{itemize}
	\itemsep -0.2em
	\cveventitem{Only open-source implementation with reproducablity}
	\cveventitem{Improving mutation and replacement modules of Genetics optimizer}
	\cveventitem{Outperforming the original paper}
\end{itemize}
\vspace{12pt}


%
\cvevent{2018}{Paper Implementation}{Open Source}{Implementation of \emph{Deep Context-aware Descreening and Rescreening of Halftone Images (Kim, T.H et al. 2018)} using PyTorch. This was the my first project using PyTorch where I learned this framework.\\
You can find the source code here: \href{https://github.com/Nikronic/Deep-Halftoning}{\underline{github.com/Nikronic/Deep-Halftoning}}}

Main contributions:
\begin{itemize}
	\itemsep -0.2em
	\cveventitem{Only open source implementation}
	\cveventitem{The implementation can be divided into below separate projects:}
	\begin{itemize}
		\itemsep -0.2em
		\cveventitem{\href{https://github.com/Nikronic/CoarseNet}{\underline{CoarseNet}}: Modified U-Net as a low-pass filter to remove halftone patterns}
		\cveventitem{\href{https://github.com/Nikronic/DetailsNet}{\underline{DetailsNet}}: A CNN generator and two discriminators for adding details}
		\cveventitem{\href{https://github.com/Nikronic/EdgeNet}{\underline{EdgeNet}}: A simple CNN to extract Canny edge features to preserve details}
		\cveventitem{\href{https://github.com/Nikronic/Halftoning-Algorithms}{\underline{Halftoning Algorithms}}: Implementation of some of the halftone algorithms provided in most recent digital color halftoning books as ground truth}
	\end{itemize}
\end{itemize}
\vspace{12pt}



%---------------------------------------------------------------------------------------
%	LANGUAGE SECTION
%--------------------------------------------------------------------------------------
\cvsection{Languages}

\begin{itemize}
	\cveventitem{\makebox[6cm]{Persian/Farsi:\hfill} Native}
	\cveventitem{\makebox[6cm]{English:\hfill} Proficient}
\end{itemize}
\vspace{12pt}


%---------------------------------------------------------------------------------------
%	REFERENCES SECTION
%--------------------------------------------------------------------------------------
\cvsection{References}\\
Available upon request.
\vspace{12pt}


%-------------------------------------------------------------------------------------------------
%	ARTIFICIAL FOOTER (fancy footer cannot exceed linewidth) 
%--------------------------------------------------------------------------------------------------

\null
\vspace*{\fill}
\hspace{-0.25\linewidth}\colorbox{bgcol}{\makebox[1.5\linewidth][c]{\mystrut \small \textcolor{white}{\href{https://nikronic.github.io}{nikronic.github.io}} $\cdot$ \textcolor{white}{\href{https://github.com/Nikronic}{github.com/Nikronic}} $\cdot$ \textcolor{white}{updated: 12.12.2021}}}

%============================================================================%
%
%
%
%	DOCUMENT END
%
%
%
%============================================================================%
\end{document}
